(1) 记$\boldsymbol d=(vk_x,vk_y,m)$，$E_\pm=\pm|\boldsymbol d|=\pm\sqrt{v^2k^2+m^2}$。下带本征矢的Bloch球自旋反平行于$\boldsymbol d$，按照题设$\mathcal A=i\langle u|\nabla u\rangle$约定，下带曲率为$\Omega_-(\boldsymbol k)=\tfrac12\hat{\boldsymbol d}\cdot(\partial_{k_x}\hat{\boldsymbol d}\times\partial_{k_y}\hat{\boldsymbol d})=mv^2/[2(m^2+v^2k^2)^{3/2}]$。它在$k=0$附近最强，随$k^{-3}$衰减，质量变号时整体反号。
(2) 定义$\mathcal C=(2\pi)^{-1}\int_{\mathbb R^2}\Omega_-\,d^2k$。极坐标给$\mathcal C=\int_0^\infty [mv^2k/(2(m^2+v^2k^2)^{3/2})]\,dk=\operatorname{sgn}(m)/2$。这不是允许晶格绝缘体的整数Chern数：常质量连续模型在$k\to\infty$方向依赖角度，无法将整个动量无穷远视为同一个紧化点，故半整数反映未完备的紫外正则化。
(3) 对$M(k)=m-Bk^2$，$\hat{\boldsymbol d}$在$k=0$指向$\operatorname{sgn}(m)\hat z$，无穷远指向$-\operatorname{sgn}(B)\hat z$，从而动量平面可紧化为球。对径向纹理，设$\cos\theta(k)=M(k)/\sqrt{v^2k^2+M(k)^2}$，方位角为$\varphi$，下带曲率积分化为$\mathcal C=\frac12[\cos\theta(0)-\cos\theta(\infty)]=[\operatorname{sgn}(m)+\operatorname{sgn}(B)]/2$，结果必为$0,\pm1$。只有$m=0$或需要跨越模型正则化参数$B=0$的间隙/边界改变时才可变换拓扑数。
(4) 对完全填满的孤立带，Hall响应与积分Berry曲率成正比，但$\sigma_{xy}$的总符号取决于电荷符号和所选$\mathcal A$、响应定义的约定，应统一固定后比较。局域规变$|u_-\rangle\to e^{i\chi(\boldsymbol k)}|u_-\rangle$改变$\mathcal A$为$\mathcal A-\nabla\chi$，曲率及$\mathcal C$不变。真实晶格的Brillouin区本身紧致，额外高能带负责完整的整数化，不能把单个无正则化Dirac锥的半整数直接称为独立可测的整数Hall平台。
